\chapter{Evaluation}
\label{ch:evaluation}

This chapter defines how forecasts are scored, how models are compared, and what ``better'' means quantitatively.

%% ──────────────────────────────────────────────
\section{Metrics}
\label{sec:eval-metrics}
%% ──────────────────────────────────────────────

Table~\ref{tab:eval-metrics} lists every metric used in the project and its role.

\begin{table}[htbp]
\centering
\caption{Forecast evaluation metrics.}
\label{tab:eval-metrics}
\small
\begin{tabular}{@{}lp{5.5cm}l@{}}
\toprule
\textbf{Metric} & \textbf{What It Measures} & \textbf{Role} \\
\midrule
$\QLIKE$                 & Quasi-likelihood loss; penalizes relative forecast errors  & Primary \\
MSE                      & Mean squared error on $\log$-$\RV$                        & Secondary \\
MAE                      & Mean absolute error on $\log$-$\RV$                       & Robustness check \\
Diebold--Mariano test    & Statistical significance of pairwise forecast differences  & ``Is model A better than B?'' \\
Model Confidence Set     & Set of models not significantly worse than the best        & ``Which models are top tier?'' \\
\bottomrule
\end{tabular}
\end{table}

The primary loss function is $\QLIKE$:
\begin{equation}
\label{eq:qlike}
\QLIKE = \frac{1}{T} \sum_{t=1}^{T} \left( \frac{\hat{\sigma}_t^2}{\sigma_t^2} - \log \frac{\hat{\sigma}_t^2}{\sigma_t^2} - 1 \right)
\end{equation}
$\QLIKE$ penalizes relative errors, not absolute ones.
It ranks forecasters consistently even when $\RV$ is measured with noise, because it belongs to the class of loss functions that are robust to imperfect volatility proxies \citep{Patton2011}.

\begin{keyidea}[$\QLIKE$ Is the Only Loss That Matters for Ranking]
$\QLIKE$ is the primary metric for all model selection and comparison decisions.
MSE and MAE serve only as secondary diagnostics.
The Diebold--Mariano test \citep{DieboldMariano1995} confirms whether a $\QLIKE$ difference is statistically significant; the Model Confidence Set \citep{HansenLundeNason2011} identifies which models belong to the top tier.
\end{keyidea}


%% ──────────────────────────────────────────────
\section{Validation Protocol}
\label{sec:validation}
%% ──────────────────────────────────────────────


%% ─── 13.2.1  Purged k-fold CV ───
\subsection{Purged k-Fold CV with Embargo}
\label{sec:purged-cv}

Standard cross-validation leaks information in time series because observations near fold boundaries share overlapping feature windows.
Purged CV removes observations within the purge window on both sides of the train/test boundary.
An embargo adds a further gap after each test fold to prevent label leakage from multi-day targets \citep{LopezdePrado2018}.

\begin{center}
\begin{tikzpicture}[
  xscale=1.0,
  >=Stealth,
]
  % Time axis
  \draw[arrow] (0,0) -- (14,0) node[right] {time};

  % Train block
  \fill[blue!25] (0.5,0.3) rectangle (5.0,1.0);
  \node[font=\small] at (2.75,0.65) {Train};

  % Purge zone
  \fill[pattern=north east lines, pattern color=red!60] (5.0,0.3) rectangle (6.0,1.0);
  \node[font=\scriptsize, red!70!black, above] at (5.5,1.0) {Purge};

  % Embargo zone
  \fill[pattern=north west lines, pattern color=orange!70] (6.0,0.3) rectangle (7.2,1.0);
  \node[font=\scriptsize, orange!70!black, above] at (6.6,1.0) {Embargo};

  % Test block
  \fill[green!25] (7.2,0.3) rectangle (11.5,1.0);
  \node[font=\small] at (9.35,0.65) {Test};

  % Purge zone (post-test)
  \fill[pattern=north east lines, pattern color=red!60] (11.5,0.3) rectangle (12.5,1.0);
  \node[font=\scriptsize, red!70!black, above] at (12.0,1.0) {Purge};

  % Brace annotations
  \draw[decorate, decoration={brace, amplitude=4pt, mirror}] (5.0,-0.1) -- (7.2,-0.1)
    node[midway, below=5pt, font=\scriptsize] {no data used};

\end{tikzpicture}
\end{center}

We use purged 5-fold CV for hyperparameter tuning.
Purge window: equal to the longest feature lookback (22 trading days for the monthly HAR component).
Embargo window: equal to the forecast horizon (1 or 5 days).


%% ─── 13.2.2  Walk-forward ───
\subsection{Walk-Forward Evaluation}
\label{sec:walk-forward}

Walk-forward is the primary out-of-sample evaluation method.
Train on a rolling 5-year window, forecast the next period, step forward by one period, and repeat.
All reported $\QLIKE$ numbers come from this procedure, never from in-sample or CV scores.

\begin{center}
\begin{tikzpicture}[xscale=0.65, >=Stealth]
  % Iteration 1
  \fill[blue!20] (0,2.4) rectangle (10,3.0);
  \node[font=\scriptsize] at (5,2.7) {Train (5 yr)};
  \fill[green!30] (10,2.4) rectangle (11,3.0);
  \node[font=\scriptsize] at (10.5,2.7) {Test};

  % Iteration 2
  \fill[blue!20] (1,1.2) rectangle (11,1.8);
  \node[font=\scriptsize] at (6,1.5) {Train (5 yr)};
  \fill[green!30] (11,1.2) rectangle (12,1.8);
  \node[font=\scriptsize] at (11.5,1.5) {Test};

  % Iteration 3
  \fill[blue!20] (2,0.0) rectangle (12,0.6);
  \node[font=\scriptsize] at (7,0.3) {Train (5 yr)};
  \fill[green!30] (12,0.0) rectangle (13,0.6);
  \node[font=\scriptsize] at (12.5,0.3) {Test};

  % Arrow showing roll
  \draw[arrow, gray] (13.5,3.0) -- (13.5,0.3) node[midway, right, font=\scriptsize] {roll};

  % Time axis
  \draw[arrow] (0,-0.5) -- (14,-0.5) node[right, font=\scriptsize] {time};
\end{tikzpicture}
\end{center}


%% ──────────────────────────────────────────────
\section{Success Target}
\label{sec:success-target}
%% ──────────────────────────────────────────────

Three criteria define success:

\begin{enumerate}[nosep]
  \item \textbf{$\QLIKE$ improvement:} 30--80 bps over the HARQ baseline \citep{BPQ2016}, averaged across the 35-instrument universe.
  \item \textbf{Statistical significance:} Diebold--Mariano $p < 0.05$ vs.\ each baseline (HAR, HARQ, SHAR, Realized GARCH).
  \item \textbf{Economic value:} Positive out-of-sample utility gain in a volatility-targeting portfolio \citep{MoreiraMuir2017}.
\end{enumerate}


%% ──────────────────────────────────────────────
%% Warning box
%% ──────────────────────────────────────────────

\begin{warning}[Train with $\QLIKE$, Not MSE]
MSE is dominated by extreme volatility days.
A single crisis observation can flip MSE rankings.
$\QLIKE$ penalizes relative forecast errors and ranks models consistently regardless of the volatility regime \citep{Patton2011}.

LightGBM does not provide a built-in $\QLIKE$ loss.
Training requires a custom objective that returns both the gradient and Hessian of the $\QLIKE$ loss with respect to the predicted value (see Chapter~\ref{ch:lightgbm}, Section~\ref{sec:qlike-objective}).
If you train with MSE and evaluate with $\QLIKE$, the model optimizes the wrong surface and rankings will not transfer.
\end{warning}


%% ──────────────────────────────────────────────
\section{Retransformation Bias}
\label{sec:retransformation}
%% ──────────────────────────────────────────────

Models forecast in $\log$-$\RV$ space, but $\QLIKE$ and downstream applications require level-space forecasts.

Naive exponentiation is biased low by Jensen's inequality because $\E[\exp(X)] > \exp(\E[X])$ for any non-degenerate random variable.
The correction is:
\begin{equation}
\label{eq:retransformation}
\widehat{\RV}_{t+1} = \exp\!\left(\widehat{\log \RV}_{t+1} + \hat{\sigma}^2_\varepsilon / 2\right)
\end{equation}
where $\hat{\sigma}^2_\varepsilon$ is the variance of log-space forecast errors, estimated from a rolling 60-day OOS window.

\begin{table}[htbp]
\centering
\caption{Approximate retransformation bias by forecast horizon.}
\label{tab:retrans-bias}
\small
\begin{tabular}{@{}lc@{}}
\toprule
\textbf{Horizon} & \textbf{Bias (without correction)} \\
\midrule
$h = 1$  & $\sim$4\%  \\
$h = 5$  & $\sim$10\% \\
$h = 22$ & $\sim$19\% \\
\bottomrule
\end{tabular}
\end{table}

\begin{keyidea}[Apply Before Level-Space $\QLIKE$]
Apply the correction before computing level-space $\QLIKE$.
Without it, every forecast is systematically low, the bias grows with horizon, and the MZ regression (Section~\ref{sec:mz}) will show $a > 0$.
\end{keyidea}


%% ──────────────────────────────────────────────
\section{Mincer--Zarnowitz Regression}
\label{sec:mz}
%% ──────────────────────────────────────────────

$\QLIKE$ tells you which model wins; the MZ regression tells you \emph{why} a forecast is bad.

\begin{equation}
\label{eq:mz}
\sigma^2_t = a + b \cdot h_t + \varepsilon_t
\end{equation}
where $\sigma^2_t$ is realized variance and $h_t$ is the forecast.
An efficient forecast satisfies $a = 0$, $b = 1$.

\begin{table}[htbp]
\centering
\caption{MZ regression diagnostics.}
\label{tab:mz-diagnostics}
\small
\begin{tabular}{@{}llp{5cm}@{}}
\toprule
\textbf{Pattern} & \textbf{Diagnosis} & \textbf{Fix} \\
\midrule
$a > 0,\; b \approx 1$ & Systematic under-prediction    & Check retransformation bias \\
$a \approx 0,\; b < 1$ & Forecast too smooth            & Increase reactivity to recent $\RV$ \\
$a \approx 0,\; b > 1$ & Forecast too volatile           & Regularize or increase shrinkage \\
$a = 0,\; b = 1$       & Efficient forecast             & Report $R^2$ \\
\bottomrule
\end{tabular}
\end{table}

Test the joint hypothesis $H_0\colon a = 0,\; b = 1$ with an $F$-test.

\begin{warning}[Use HAC Standard Errors]
Use Newey--West (HAC) standard errors.
Vol forecast errors are serially correlated; OLS standard errors are too small and reject $H_0$ too often.
\end{warning}


%% ──────────────────────────────────────────────
\section{Diebold--Mariano Test}
\label{sec:dm-test}
%% ──────────────────────────────────────────────

A $\QLIKE$ improvement means nothing without a $p$-value.

Define the loss differential:
\begin{equation}
\label{eq:dm-loss-diff}
d_t = L(\sigma^2_t,\, h^A_t) - L(\sigma^2_t,\, h^B_t)
\end{equation}

The test statistic is:
\begin{equation}
\label{eq:dm-stat}
\mathrm{DM} = \frac{\bar{d}}{\sqrt{\widehat{\mathrm{Var}}(\bar{d})}}
\end{equation}
with Newey--West HAC variance estimator, bandwidth $\ell = \lfloor T^{1/3} \rfloor$.
Under $H_0$ of equal predictive ability, $\mathrm{DM} \xrightarrow{d} \N(0,1)$.
Reject if $|\mathrm{DM}| > 1.96$ at 5\%.

\begin{keyidea}[Every $\QLIKE$ Comparison Needs a DM $p$-Value]
Run pairwise: ML vs HAR, ML vs HARQ, ML vs SHAR, ML vs Realized GARCH.
If $p > 0.05$, the improvement is not credible.
\end{keyidea}

\begin{warning}[Small-Sample Correction]
With $T < 100$, use the Harvey--Leybourne--Newbold correction \citep{HarveyLeybourneNewbold1997}: replace $\N(0,1)$ with $t_{T-1}$ and apply the finite-sample correction factor.
\end{warning}


%% ──────────────────────────────────────────────
\section{Model Confidence Set}
\label{sec:mcs}
%% ──────────────────────────────────────────────

The DM test compares two models.
The MCS compares all of them simultaneously without inflating the false-positive rate.

\textbf{Procedure:}
\begin{enumerate}[nosep]
  \item Start with the full model set $\mathcal{M}_0$.
  \item Test $H_0$: all models in the current set have equal expected loss.
  \item If rejected, remove the worst model (highest average loss).
  \item Repeat until $H_0$ is not rejected; survivors form $\widehat{\mathcal{M}}^*_\alpha$.
\end{enumerate}

The MCS $p$-value for each model is the smallest $\alpha$ at which the model is excluded.
Report it for every model.

\begin{table}[htbp]
\centering
\caption{MCS reporting template.}
\label{tab:mcs-template}
\small
\begin{tabular}{@{}lcccc@{}}
\toprule
\textbf{Model} & $\QLIKE$ & \textbf{DM vs HARQ} & \textbf{MCS $p$} & \textbf{In MCS$_{90\%}$?} \\
\midrule
LightGBM (L0--2) & ---   & ---      & ---   & --- \\
HARQ              & ---   & baseline & 1.000 & Yes \\
HAR               & ---   & ---      & ---   & --- \\
SHAR              & ---   & ---      & ---   & --- \\
Realized GARCH    & ---   & ---      & ---   & --- \\
\bottomrule
\end{tabular}
\end{table}

\begin{keyidea}[The Credibility Test]
The MCS is the credibility test for the presentation.
If LightGBM and plain HAR both survive in the 90\% MCS, you cannot claim ML superiority.
If HAR is excluded and LightGBM survives, that is defensible.
Use \texttt{arch.bootstrap.MCS} in Python.
\end{keyidea}

\begin{figure}[htbp]
\centering
\begin{tikzpicture}[node distance=1.4cm]
  % Step 1
  \node[flowblock, minimum width=6cm] (start)
    {Start: all $M$ candidate models};

  % Step 2
  \node[flowblock, minimum width=6cm, below=of start] (test)
    {Test $H_0$: equal predictive ability};

  % Decision
  \node[decisionblock, below=of test] (decision)
    {$H_0$ rejected?};

  % Remove worst
  \node[flowblock, minimum width=4cm, right=2.5cm of decision] (remove)
    {Remove worst model};

  % Survivors
  \node[flowblock, minimum width=6cm, below=1.6cm of decision] (survivors)
    {Survivors $= \widehat{\mathcal{M}}^*_\alpha$};

  % Arrows
  \draw[arrow] (start) -- (test);
  \draw[arrow] (test) -- (decision);
  \draw[arrow] (decision) -- node[left, font=\scriptsize] {No} (survivors);
  \draw[arrow] (decision) -- node[above, font=\scriptsize] {Yes} (remove);
  \draw[arrow] (remove) |- (test);
\end{tikzpicture}
\caption{MCS elimination procedure.}
\label{fig:mcs-flowchart}
\end{figure}


%% ──────────────────────────────────────────────
\section{Deflated Sharpe Ratio}
\label{sec:dsr}
%% ──────────────────────────────────────────────

If the forecast feeds a trading strategy, the backtest Sharpe must survive correction for the number of configurations tested.

Testing $N$ strategies inflates the expected best Sharpe.
Under pure luck, $\E[\max_i \SR_i] \approx \sqrt{2 \ln N}$; for $N = 20$ this is $\approx 2.45$.

\begin{equation}
\label{eq:dsr}
\DSR = \Phi\!\left(\frac{(\widehat{\SR} - \SR_0)\sqrt{T-1}}{\sqrt{1 - \hat{\gamma}_3 \widehat{\SR} + \frac{\hat{\gamma}_4 - 1}{4}\widehat{\SR}^2}}\right)
\end{equation}
where $\widehat{\SR}$ is the observed annualized Sharpe, $\SR_0 = \sqrt{2\ln N}$ is the luck threshold, $T$ is the number of return observations, $\hat{\gamma}_3$ is skewness, and $\hat{\gamma}_4$ is kurtosis.

Decision: $\DSR > 0.95$ is credible.

\begin{warning}[Every Experiment Counts]
Every experiment counts as a trial.
Every hyperparameter grid point, every feature set, every quick look increments $N$.
Log experiments from day one.
\end{warning}


%% ──────────────────────────────────────────────
\section{Evaluation Pitfalls}
\label{sec:eval-pitfalls}
%% ──────────────────────────────────────────────

Six errors that invalidate results.

\begin{warning}[What Invalidates Your Results]
\begin{enumerate}[nosep]
  \item \textbf{Random $k$-fold on time series.} Trains on the future.
    Always purged CV or walk-forward.
  \item \textbf{$\QLIKE$ improvement without DM test.} A 3\% gain with $p = 0.12$ is noise.
  \item \textbf{Sharpe without DSR.} A Sharpe of 1.5 from 20 experiments is below the luck threshold.
  \item \textbf{Training on one regime, testing on another.} 2015--2019 train, 2020 test is a stress test, not general evaluation.
  \item \textbf{Lookahead in features.} Day-$t$ VIX close for day-$t$ $\RV$ is look-ahead.
    All features must precede the forecast cutoff (Chapter~\ref{ch:pipeline}).
  \item \textbf{Tiny improvements without economic significance.} 0.5\% $\QLIKE$ gain is statistically real but economically meaningless after costs.
\end{enumerate}
\end{warning}


%% ──────────────────────────────────────────────
\section{Evaluation Workflow}
\label{sec:eval-workflow}
%% ──────────────────────────────────────────────

Every forecasting experiment follows this pipeline.

\begin{figure}[htbp]
\centering
\begin{tikzpicture}[node distance=1.2cm]
  % Nodes
  \node[flowblock, minimum width=7cm, fill=prereqpurple!10] (holdout)
    {Reserve holdout (final 6 months)};

  \node[flowblock, minimum width=7cm, fill=blue!8, below=of holdout] (initlog)
    {Initialize experiment log ($N = 0$)};

  \node[flowblock, minimum width=7cm, fill=blue!8, below=of initlog] (tune)
    {Tune: purged $k$-fold CV};

  \node[flowblock, minimum width=7cm, fill=blue!8, below=of tune] (walkfwd)
    {Walk-forward OOS: $\QLIKE$, MSE};

  \node[flowblock, minimum width=7cm, fill=blue!8, below=of walkfwd] (mzreg)
    {MZ regression: $a = 0$, $b = 1$?};

  \node[flowblock, minimum width=7cm, fill=blue!8, below=of mzreg] (dmtest)
    {DM test: pairwise vs baselines};

  \node[flowblock, minimum width=7cm, fill=keyorange!10, below=of dmtest] (mcs)
    {MCS: top-tier model set};

  \node[decisionblock, below=1.2cm of mcs] (strategy)
    {Strategy?};

  \node[flowblock, minimum width=4cm, fill=memgold!15, right=2.5cm of strategy] (dsr)
    {DSR on Sharpe};

  \node[flowblock, minimum width=7cm, fill=intgreen!10, below=1.6cm of strategy] (report)
    {Report all metrics};

  % Side annotation
  \node[font=\scriptsize, text=gray, right=0.3cm of tune] {Every experiment increments $N$};

  % Arrows — main spine
  \draw[arrow] (holdout) -- (initlog);
  \draw[arrow] (initlog) -- (tune);
  \draw[arrow] (tune) -- (walkfwd);
  \draw[arrow] (walkfwd) -- (mzreg);
  \draw[arrow] (mzreg) -- (dmtest);
  \draw[arrow] (dmtest) -- (mcs);
  \draw[arrow] (mcs) -- (strategy);

  % Decision branches
  \draw[arrow] (strategy) -- node[above, font=\scriptsize] {Yes} (dsr);
  \draw[arrow] (strategy) -- node[left, font=\scriptsize] {No} (report);
  \draw[arrow] (dsr) |- (report);
\end{tikzpicture}
\caption{Evaluation workflow. Every forecasting experiment follows this pipeline end to end.}
\label{fig:eval-workflow}
\end{figure}

\begin{keyidea}[No Guarantees, Only Evidence]
This workflow does not guarantee a good forecast.
It guarantees that if you find one, the evidence survives review.
\end{keyidea}

\begin{table}[htbp]
\centering
\caption{Evaluation toolkit summary.}
\label{tab:eval-toolkit}
\small
\begin{tabular}{@{}llc@{}}
\toprule
\textbf{Tool} & \textbf{Question It Answers} & \textbf{Reference} \\
\midrule
$\QLIKE$            & Which forecast has lower loss?       & \citet{Patton2011} \\
Retransformation    & Is my level-space forecast biased?   & \citet{Patton2011} \\
Mincer--Zarnowitz   & Is the forecast calibrated?          & \citet{MincerZarnowitz1969} \\
Diebold--Mariano    & Is the improvement significant?      & \citet{DieboldMariano1995} \\
Model Confidence Set & Which models are top tier?           & \citet{HansenLundeNason2011} \\
Purged $k$-fold CV  & Am I leaking future info?            & \citet{LopezdePrado2018} \\
Deflated Sharpe     & Is the backtest Sharpe real?          & \citet{Bailey2014DSR} \\
\bottomrule
\end{tabular}
\end{table}
